\pagebreak
    \hspace{0pt}
    \vfill
    \addcontentsline{toc}{section}{Abstract}
    \section*{Abstract}
    
    Conventional neuro-oncological evaluation in Glioblastoma (GBM) care is primarily depending on static volumetric measurements and standardized response criteria to characterize prior tumor changes rather than forecast future tumor behavior. This study introduces the \texttt{NEXUS Brain Virtual Human Twin (VHT)} framework to bridge retrospective imaging analysis and prospective clinical decision making. The framework integrates longitudinal \texttt{NIfTI} segmentations, treatment exposure, and molecular context into a mechanistic, patient specific model of tumor growth and treatment response. An automated pipeline converted clinical and imaging repositories into structured digital patient records, retaining 155 of 203 patients (76.3\%) and yielding 546 longitudinal observations. A Hierarchical Bayesian Model (HBM), using non-centered parameterization and No-U-Turn Sampler (NUTS) \texttt{MCMC}, was formulated from the Fisher-KPP reaction diffusion equation, \(v_i=2\sqrt{D_i\rho_i}\). It simultaneously estimated population level hyperpriors and patient specific proliferation, diffusion, and carrying capacity parameters \((\rho_i,D_i,K_i)\), together with treatment efficacy coefficients for Chemotherapy, Radiotherapy, Immunotherapy, additional systemic therapies, and Brachytherapy. Calibrated parameters were deployed in a forecasting engine and interactive \texttt{Streamlit} dashboard for counterfactual scenario exploration. Cox proportional hazards models evaluated image derived kinetic features alongside \textit{IDH1}, \textit{MGMT} methylation, 1p/19q codeletion, age and etc. Convergence diagnostics showed \(\widehat{R}\leq1.04\), zero divergences, and \(\sigma_{\mathrm{obs},\log v}=0.8720\) (94\% HPD: 0.8190--0.9260). Patient specific parameters exhibited substantial heterogeneity with median proliferation was \(0.067\text{ day}^{-1}\), median diffusion \(0.434\text{ mm}^2/\text{day}\), and median wavefront speed \(0.086\text{ mm/day}\). Counterfactual simulations projected prospective tumor trajectories and evaluated multimodal interventions, including virtual resection and combined therapy windows, yielding substantial volume reductions. The VHT-driven Cox model produced a significant likelihood ratio test \((p=4.31\times10^{-7})\), identifying enhancing tumor velocity as a significant prognostic factor where one standard deviation increase was associated with an approximately 35.4\% higher mortality hazard. \textit{IDH1} mutation was protective. The \textsc{Nexus Brain} VHT framework translates longitudinal imaging into actionable growth projections and counterfactual simulations, demonstrating that dynamic image derived kinetics provide prognostic information beyond static baseline volumes and establishing a patient specific foundation for precision
    
    \vfill
    \hspace{0pt}
\pagebreak